\makeabstract
{
Nonlinear Independent Component Analysis can identify spatial structure from the neural activity of hippocampal CA1 neurons
}
{
Ichika Abe (1), Takatomi Kubo, PhD (1)
}
{
\textsuperscript{1}Graduate School of Science and Technology, Nara Institute of Science and Technology, Ikoma, Japan
}
{
The purpose of this study is to investigate whether spatial structures can be extracted from high-dimensional neural recordings. Previous work has shown that populations of neurons, including place cells and boundary vector cells, encode spatial information. However, neural activity is often high-dimensional and nonlinear, making it unclear whether and how spatial structure can emerge in low-dimensional representations. To extract latent structures from high-dimensional neural activity, we applied a nonlinear Independent Component Analysis (ICA) approach to neural data collected during spatial navigation tasks. In particular, we adopted CEBRA (Consistent EmBeddings of high-dimensional Recordings using Auxiliary variables), a recent contrastive learning framework that leverages auxiliary variables.
}
{
The experiment consisted of two phases. In the preliminary phase, we used a synthetic neural dataset to introduce complementary features and enhance data diversity. In the main phase, we analysed hippocampal Cornu Ammonis area 1 (CA1) recordings from freely moving rats to examine whether the insights gained from the synthetic data generalize to real neural dynamics. CEBRA was applied in both discovery-driven (self-supervised) and hypothesis-driven (supervised) modes to extract low-dimensional, behaviourally relevant latent representations from the neural activity.
}
{
In the preliminary experiments, the discovery-driven mode produced latent representations whose structures closely matched the geometry of the corresponding environments. To further assess structural alignment, we trained embeddings in the hypothesis-driven mode using the rats’ positional coordinates as auxiliary variables, and quantified the similarity between embedding and position using the R² score from a linear regression fit. We applied the same procedure in the main experiments. In the preliminary experiments, the discovery-driven mode achieved a higher R² score (\textasciitilde{}0.94) than the hypothesis-driven mode, indicating that spatial structure could be directly extracted from the autocorrelation of neural data. In contrast, in the main experiments, the hypothesis-driven mode achieved superior performance (\textasciitilde{}0.91), suggesting that conditioned non-stationarity in neural activity can more effectively reveal spatially relevant structure. 
}
{
Our findings reveal that rich spatial features are embedded in neural activity and can be recovered through suitable low-dimensional embeddings. The spatial structure observed in the learned latent space suggests that neural population activity may be organized along a low-dimensional manifold reflecting the geometry of physical space.
}

\newpage
